\subsection{Почему не используют матричное решение}

Заметим, что для получения ответа нам нужно обратить матрицу $X^T X$.
Это создаёт множество проблем:

\begin{enumerate}
    \item Основная проблема в обращении матрицы -- это то, что вычислительно обращать большие матрицы дело
    сложное, а мы бы хотели работать с датасетами, в которых у нас могут быть миллионы точек,
    \item Матрица $X^T X$, хотя почти всегда обратима в разумных задачах машинного обучения, часто плохо обусловлена.
    Особенно если признаков много, между ними может появляться приближённая линейная зависимость, которую мы
    можем упустить на этапе формулировки задачи.
    Тогда погрешность нахождения $w$ будет зависеть от квадрата числа обусловленности матрицы $X$, что очень плохо.
    Это делает полученное таким образом решение численно неустойчивым: малые возмущения $y$ могут приводить к
    катастрофическим изменениям $w$.
\end{enumerate}

Полностью вылечить проблемы мы не сможем, но никто и не обязывает нас останавливаться на «точном» решении (
которое всё равно никогда не будет вполне точным).
Поэтому используют другой метод

Минимизируемый функционал является гладким и выпуклым, а это значит, что можно эффективно искать точку его
минимума с помощью итеративных градиентных методов.
Далее об одном самом базовом подходе.

Как известно, градиент функции в точке направлен в сторону её наискорейшего роста, а антиградиент (противоположный
градиенту вектор) в сторону наискорейшего убывания.
То есть имея какое-то приближение оптимального значения параметра $w$, мы можем его улучшить, посчитав градиент
функции потерь в точке и немного сдвинув вектор весов в направлении антиградиента:
\[ w_j \rightarrow w_j - \alpha \frac{d}{dw_j} L(f_w,X,y) \],
где $\alpha$ -- это параметр алгоритма («темп обучения»), который контролирует величину шага в направлении
антиградиента.
Описанный алгоритм называется градиентным спуском.

Посмотрим, как будет выглядеть градиентный спуск для функции потерь $L(f_w,X,y) = \frac{1}{N} \abs{X w - y}_2^2$.
Градиент квадрата евклидовой нормы мы уже считали; соответственно,

\[  \nabla_w L = \frac{2}{N} X^T (X w - y) \]

Следовательно, стартовав из какого-то начального приближения, мы можем итеративно уменьшать значение функции,
пока не сойдёмся (по крайней мере в теории) к минимуму (вообще говоря, локальному, но в данном случае глобальному)

Несколько общих замечаний:

\begin{itemize}

    \item Поскольку задача выпуклая, выбор начальной точки влияет на скорость сходимости, но не настолько
    сильно, чтобы на практике нельзя было стартовать всегда из нуля или из любой другой приятной вам точки;
    \item Число обусловленности матрицы $X$ существенно влияет на скорость сходимости градиентного спуска: чем
    более вытянуты эллипсоиды уровня функции потерь, тем хуже;
    \item Темп обучения $\alpha$ тоже сильно влияет на поведение градиентного спуска; вообще говоря, он является
    гиперпараметром алгоритма, и его, возможно, придётся подбирать отдельно.
    Другими гиперпараметрами являются максимальное число итераций $S$ и (или) порог tolerance.

\end{itemize}

Вычислительная сложность градиентного спуска -- $O(NDS)$, где, как и выше, $N$ -- длина выборки, $D$ -- число
признаков у одного объекта.
Сравните с оценкой $O(D^2 N+D^3)$ для «наивного» вычисления аналитического решения.

Сложность по памяти -- $O(ND)$ на хранение выборки.
В памяти мы держим и выборку, и градиент, но в большинстве реалистичных сценариев доминирует выборка.

\subsection{Стохастический градиентный спуск, преимущества}

На каждом шаге градиентного спуска нам требуется выполнить потенциально дорогую операцию вычисления градиента по
всей выборке (сложность $O(ND)$).
Возникает идея заменить градиент его оценкой на подвыборке (в английской литературе такую подвыборку обычно именуют
batch или mini-batch; в русской разговорной терминологии тоже часто встречается слово батч или мини-батч).

А именно, если функция потерь имеет вид суммы по отдельным парам объект-таргет
\[ L(w,X,y) = \frac{1}{N} \sum_{i=1}^N L(w,x_i,y_i)\]

а градиент, соответственно, записывается в виде
\[ \nabla_w L(w,X,y) = \frac{1}{N} \sum_{i=1}^N \nabla_w L(w,x_i,y_i) \]

то предлагается брать оценку
\[ \nabla_w L(w,X,y) = \frac{1}{B} \sum_{t=1}^B \nabla_w L(w,x_{i_t},y_{i_t}) \]
для некоторого подмножества этих пар $(x_{i_t},y_{i_t})_{t=1}^B$.

Как делить выборку на батчи?
Ясно, что можно было бы случайным образом сэмплировать их из полного датасета, но даже если использовать быстрый
алгоритм вроде резервуарного сэмплирования, сложность этой операции не самая оптимальная.
Поэтому используют линейный проход по выборке (которую перед этим лучше всё-таки случайным образом перемешать).
Давайте введём ещё один параметр нашего алгоритма: размер батча, который мы обозначим $B$.
Теперь на $B$ очередных примерах вычислим градиент и обновим веса модели.
При этом, вместо количества шагов алгоритма обычно задают количество эпох $E$.
Это ещё один гиперпараметр.
Одна эпоха -- это один полный проход нашего сэмплера по выборке.
Заметим, что если выборка очень большая, а модель компактная, то даже первый проход бывает можно не заканчивать

Сложность по времени -- $O(NDE)$.
На первый взгляд, она такая же, как и у обычного градиентного спуска, но заметим, что мы сделали в $\frac{N}{B}$
раз больше шагов, то есть веса модели претерпели намного больше обновлений.

Сложность по памяти можно довести до $O(BD)$: ведь теперь всю выборку не надо держать в памяти, а достаточно
загружать лишь текущий батч (а остальная выборка может лежать на диске, что удобно, так как в реальности задачи, в
которых выборка целиком не влезает в оперативную память, встречаются сплошь и рядом).
Заметим, впрочем, что при этом лучше бы $B$ взять побольше: ведь чтение с диска -- намного более затратная по
времени операция, чем чтение из оперативной памяти.

Шаги стохастического градиентного спуска заметно более шумные, но считать их получается значительно быстрее.
Они тоже сходятся к оптимальному значению потому, что матожидание оценки градиента на батче равно самому градиенту.
По крайней мере, сходимость можно получить при хорошо подобранных коэффициентах темпа обучения в случае выпуклого
функционала качества.

Существует определённая терминологическая путаница, иногда стохастическим градиентным спуском называют версию
алгоритма, в которой размер батча равен единице (то есть максимально шумная и быстрая версия алгоритма), а версии с
гораздо большим размером батча называют batch gradient descent.
В книгах, которые, возможно, старше вас, такая процедура иногда ещё называется incremental gradient descent.
Это не очень принципиально, но вы будьте готовы, если что